\documentclass[10pt,a4paper]{amsart}
\usepackage[margin=19mm]{geometry}
\usepackage{amsmath,amssymb,mathtools}
\usepackage[scr=rsfs,cal=euler]{mathalfa}
\usepackage{xfrac}
\usepackage[unicode,hidelinks,bookmarks=false]{hyperref}
\newcommand{\ie}{i.e.\ }
\newcommand{\cf}{cf.\ }
\newcommand{\resp}{respectively\ }
\newcommand{\Subsection}{\subsection}
\providecommand{\nobreakdash}{\nobreak\hbox{-}}
\setlength{\emergencystretch}{2em}
\multlinegap0pt
\usepackage{amssymb}
\usepackage{mathtools}
\usepackage{xcolor}
\usepackage{tikz}
\usepackage{float}
\newtheorem{lemma}{Lemma}
\newtheorem{theorem}{Theorem}
\usepackage{cleveref}
\crefname{lemma}{Lemma}{Lemmas}
\crefname{theorem}{Theorem}{Theorems}
\newcommand{\hr}[0]{\Tilde{\alpha}}
\newcommand{\newmu}[0]{\alpha_I}
\newcommand{\qpartition}[1]{\wp_q\left(#1\right)}
\newcommand{\speclin}[0]{\mathfrak{sl}_{r+1}(\mathbb{C})}
\newcommand{\thealtset}[0]{\mathcal{A}(\hr,\newmu)}
\title{On the $q$-multiplicity of sums of distinct simple roots of $\mathfrak{sl}_{r+1}(\mathbb{C})$}
\author{Matt McClinton}
\date{Source: arXiv:2603.20053v1 (CC BY 4.0)}
\begin{document}
\maketitle

\noindent\emph{Source and licence.} On the $q$-multiplicity of sums of distinct simple roots of $\mathfrak{sl}_{r+1}(\mathbb{C})$. Version arXiv:2603.20053v1 (CC BY 4.0), distributed by arXiv under the Creative Commons Attribution 4.0 International licence, http://creativecommons.org/licenses/by/4.0/, as recorded in the arXiv licence metadata for this exact version and shown on the abstract page https://arxiv.org/abs/2603.20053v1. Full licence notice: \texttt{licenses/arxiv-2603.20053-CC-BY-4.0.txt}.\par\smallskip

\noindent\textit{Editorial scope.} Counted unit: the article's lemma product (source0.tex lines 700--707). The article's complete original proof of it is delivered with it (source0.tex lines 708--764, one \texttt{proof} environment of 57 lines). No mathematics is written, repaired or re-derived: the statement and the proof are the article's own lines, in the article's own notation, and no retained line is substituted or re-flowed.  Retained uncounted as the source-local context the proof needs: lines 223--233; lines 300--305; lines 392--397.  Nothing was omitted from any retained span, so the package carries no omission record.  No retained line contains a citation, so the wrapper carries no bibliography and no bibliography entry.  No retained line contains a figure environment, an image-inclusion command or drawing code, so no image is delivered and the package declares no asset dependency.  Editorial formatting only, in addition: an AMS \texttt{amsart} preamble that restates the article's own declarations and macros used by the retained lines, namely the environments \texttt{lemma}, \texttt{theorem} on separate counters, together with the standard packages those lines need.  The retained environments print in this document as Lemma 1, Theorem 1 in order of appearance, whereas the article prints the same items with the numbering of its own full document.  The complete native source of the article as distributed in the arXiv e-print for this exact version is bundled unchanged as \texttt{sources/196823/source0.tex}.
\begin{lemma}\label{def:n(I^c)}
Fix a nonempty indexing set $I\subseteq[r]$.
Then the number of intervals in the interval partition of $I^c$ is given by 
\begin{align}
    n(I^c)=\begin{cases}
        n(I)-1 & \mbox{if $1,r\in I$} \\
        n(I) & \mbox{if $1\notin I$ and $r\in I$, or if $1\notin I$ and $r\in I$}\\
        n(I)+1 & \mbox{if $1,r\notin I$}.
    \end{cases}
\end{align}
\end{lemma}

\begin{theorem}\label{thm:whats in the alt set} 
Let $\hr$ be the highest root in $\speclin$.
    Let $I\subseteq[r]$ be nonempty, and let $\mathcal{I}$ be a (possibly empty) collection of nonconsecutive integers in $I^c\setminus\{1,r\}$.     
    Then $\sigma\in\thealtset$ if and only if 
    \[\sigma=1 \quad \text{or} \quad \sigma=\prod_{i\in \mathcal{I}}s_i.\]
\end{theorem}

\begin{lemma}\label{lemma: split up partition function}
    Let $I\subseteq [r]$ be any nonempty indexing set with interval partition $I=\bigsqcup_{x=1}^{n(I)}[i_x,j_x]$. Then,
    \begin{equation}
        \wp_q(\alpha_I) = \prod_{x=1}^{n(I)}\wp_q(\alpha_{i_x,j_x}).\notag
    \end{equation}
\end{lemma}

\begin{lemma}\label{lemma: mu a sum of positive roots turns KWMF into a product}
    Let $\hr$ denote the highest root in $\speclin$. Fix a nonempty indexing set $I\subseteq [r]$ with interval partition $I=\bigsqcup_{x=1}^{n(I)}[i_x,j_x]$. If $I^c=  
    \bigsqcup_{x=1}^{n(I^c)}[i'_x,j'_x]$ with $n(I^c)$ as in \Cref{def:n(I^c)}, then 
    \begin{equation}\label{eq:KWMF into a product}
        m_q(\hr_r,\alpha_I)=\prod_{x=1}^{n(I^c)}\sum_{\substack{\sigma\in\mathcal{A}(\hr,\alpha_I) \\ \sigma=\sigma_x}}(-1)^{\ell(\sigma_x)}\qpartition{\alpha_{i'_x,j'_x}-\sum_{i\in\mathcal{I}_x}\alpha_i},
    \end{equation}
    where, for each $x\in [n(I^c)]$, $\mathcal{I}_x=\{x_1,x_2,\ldots,x_m\}$ is a (possibly empty) set of nonconsecutive integers in the interval $[i'_x,j'_x]$, $\sigma_x=\prod_{i\in \mathcal{I}_x}s_i$, and we set $\sum_{i\in\mathcal{I}_x}\alpha_i=0$ if $\sigma_x=1$.  
\end{lemma}

\begin{proof}
    By \Cref{thm:whats in the alt set}, $\sigma\in\thealtset$ if and only if $\sigma$ is of the form
    \[\sigma=1 \quad \text{ and } \quad \sigma= \sigma_1\sigma_2\ldots\sigma_{n(I^c)},\]
    where, for each $x\in [n(I^c)]$, $\mathcal{I}_x$ is a (possibly empty) set of nonconsecutive integers in the interval $[i_x',j_x']$, and $\sigma_x=\prod_{i\in\mathcal{I}_x}s_i$. 
    Let $\sigma=\sigma_1\sigma_2\cdots\sigma_{n(I^c)}$ be any arbitrary element of the Weyl alternation set, and consider Kostant's partition function in the corresponding term in $m_q(\hr,\alpha_I)$:
    \begin{equation}
        \qpartition{\sigma(\hr_r+\rho_r)-\rho_r-\alpha_I} = \qpartition{\hr_r - \alpha_{I}-\sum_{x\in[n(I^c)]}\sum_{i\in\mathcal{I}_x}\alpha_i} = \qpartition{\alpha_{I^c}-\sum_{x\in[n(I^c)]}\sum_{i\in\mathcal{I}_x}\alpha_i}.\notag
    \end{equation}
    Since $I^c = \bigsqcup_{x=1}^{n(I^c)}[i'_x,j'_x]$, 
    \begin{align}
        &\qpartition{\alpha_{I^c}-\sum_{x\in[n(I^c)]}\sum_{i\in\mathcal{I}_x}\alpha_i}\notag\\
        &\qquad \qquad = \qpartition{\left(\alpha_{i'_1,j'_1}-\sum_{i\in\mathcal{I}_1}\alpha_i\right)+\left(\alpha_{i'_2,j'_2}-\sum_{i\in\mathcal{I}_2}\alpha_i\right)+\cdots+\left(\alpha_{i'_{n(I^c)},j'_{n(I^c)}}-\sum_{i\in\mathcal{I}_{n(I^c)}}\alpha_i\right)}\notag \\
        &\qquad\qquad = \qpartition{\sum_{x=1}^{n(I^c)}\left(\alpha_{i'_x,j'_x}-\sum_{i\in\mathcal{I}_x}\alpha_i\right)}.\notag
    \end{align}
    By \Cref{lemma: split up partition function},
    \begin{align}
        \qpartition{\sum_{x=1}^{n(I^c)}\left(\alpha_{i'_x,j'_x}-\sum_{i\in\mathcal{I}_x}\alpha_i\right)}
        &=\prod_{x=1}^{n(I^c)}\qpartition{\alpha_{i'_x,j'_x}-\sum_{i\in\mathcal{I}_x}\alpha_i}.\label{break into product 1}
    \end{align}
    Additionally, for any arbitrary $\sigma=\sigma_1\sigma_2\cdots\sigma_{n(I^c)}$, since $\sigma_x$ contains simple roots whose indices are disjoint from every other simple root in $\sigma_y$, then
    \begin{equation}
        \ell(\sigma) = \ell(\sigma_1\sigma_2\cdots \sigma_{n(I^c)}) = \sum_{x=1}^{n(I^c)} \ell(\sigma_x),\notag
    \end{equation}
    which implies that
    \begin{equation}
        (-1)^{\ell(\sigma)} = \prod_{x=1}^{n(I^c)}(-1)^{\ell(\sigma_x)}.\label{negative one stuff}
    \end{equation}
    Therefore, since our choice of term in $m_q(\hr,\alpha_I)$ was arbitrary, by \Cref{break into product 1} and \Cref{negative one stuff}, the weight $q$-multiplicity satisfies,
    \begin{align}
        m_q(\hr_r,\alpha_I) &= \sum_{\substack{\sigma\in\thealtset \\ \sigma=\sigma_1\sigma_2\cdots \sigma_{n(I^c)}}}\prod_{x=1}^{n(I^c)}(-1)^{\ell(\sigma_x)}\qpartition{\alpha_{i'_x,j'_x}-\sum_{i\in\mathcal{I}_x}\alpha_i}.\label{break into product 2}
    \end{align}
    Notice that \Cref{break into product 2} is a sum of terms, each of the form,
    \[(-1)^{\ell(\sigma_1)}\qpartition{\alpha_{i'_1,j'_1}-\sum_{i\in\mathcal{I}_1}\alpha_i}\cdot(-1)^{\ell(\sigma_2)}\qpartition{\alpha_{i'_2,j'_2}-\sum_{i\in\mathcal{I}_2}\alpha_i} \cdots (-1)^{\ell(\sigma_{n(I^c)})}\qpartition{\alpha_{i'_{n(I^c)},j'_{n(I^c)}}-\sum_{i\in\mathcal{I}_{n(I^c)}}\alpha_i} \]
    corresponding to some element of the Weyl alternation set $\sigma=\sigma_1\sigma_2\cdots\sigma_{n(I^c)}$. Note that for any $x\in[n(I^c)]$, we set $\sum_{i\in\mathcal{I}_x}\alpha_i=0$ if $\sigma_x=1$. To complete the proof, we show that by grouping terms together through factoring common $\sigma_x$ factors, we can state $m_q(\hr_r,\alpha_I)$ as a product of sums in the desired form. 
    
    Let $m\in[n(I^c)]$, and consider $\sigma_m=\prod_{i\in\mathcal{I}_m}s_i$ where $\mathcal{I}_m$ is a (possibly empty) fixed collection of integers in the interval $[i'_m,j'_m]$. Consider every term in \Cref{break into product 2} containing a factor corresponding to $\sigma_m$. 
    By factoring, the collection of these terms can be expressed as
    \begin{equation}
        (-1)^{\ell(\sigma_m)}\qpartition{\alpha_{i'_m,j'_m}-\sum_{i\in\mathcal{I}_m}\alpha_i}\cdot \sum_{\substack{\sigma\in\thealtset\\ \sigma=\sigma_1\sigma_2\cdots\sigma_{m-1}\sigma_{m+1}\cdots\sigma_{n(I^c)}}}\prod_{\substack{x=1\\x\neq m}}^{n(I^c)} (-1)^{\ell(\sigma_x)}\qpartition{\alpha_{i'_x,j'_x}-\sum_{i\in\mathcal{I}_x}\alpha_i}.
    \end{equation}
    Therefore, by repeating this factorization for every collection of nonconsecutive integers $\mathcal{I}_m$ in the interval $[i'_m,j'_m]$, \Cref{break into product 2} can be expressed as
    \begin{align}
        &\sum_{\substack{\sigma\in\mathcal{A}(\hr,\alpha_I) \\ \sigma=\sigma_1\sigma_2\cdots\sigma_n}}\prod_{x=1}^{n(I^c)}(-1)^{\ell(\sigma_x
        )}\qpartition{\alpha_{i'_x,j'_x}-\sum_{i\in\mathcal{I}_x}\alpha_i}\notag \\
        &= \sum_{\substack{\sigma\in\thealtset\\ \sigma=\sigma_m}}\left((-1)^{\ell(\sigma_m)}\qpartition{\alpha_{i'_m,j'_m}-\sum_{i\in\mathcal{I}_m}\alpha_i}\cdot \sum_{\substack{\sigma\in\thealtset\\ \sigma=\sigma_1\sigma_2\cdots \sigma_{m-1}\sigma_{m+1}\cdots\sigma_n}}\prod_{\substack{x=1\\x\neq m}}^{n(I^c)} (-1)^{\ell(\sigma_x)}\qpartition{\alpha_{i'_x,j'_x}-\sum_{i\in\mathcal{I}_x}\alpha_i}\right).\label{eq: KWMF with sigma 1 factored out}
    \end{align}
    In \Cref{eq: KWMF with sigma 1 factored out}, we remove the common $\sigma_m$ factor from each grouping of terms, and express the equation as
    \begin{align}
        &\sum_{\substack{\sigma\in\thealtset\\ \sigma=\sigma_m}}\left((-1)^{\ell(\sigma_m)}\qpartition{\alpha_{i'_m,j'_m}-\sum_{i\in\mathcal{I}_m}\alpha_i}\cdot \sum_{\substack{\sigma\in\thealtset\\ \sigma=\sigma_1\sigma_2\cdots \sigma_{m-1}\sigma_{m+1}\cdots\sigma_n}}\prod_{\substack{x=1\\x\neq m}}^{n(I^c)} (-1)^{\ell(\sigma_x)}\qpartition{\alpha_{i'_x,j'_x}-\sum_{i\in\mathcal{I}_x}\alpha_i}\right)\notag \\
        &\qquad\qquad=\left(\sum_{\substack{\sigma\in\thealtset\\ \sigma=\sigma_m}}(-1)^{\ell(\sigma_m)}\qpartition{\alpha_{i'_m,j'_m}-\sum_{i\in\mathcal{I}_m}\alpha_i}\right)\notag\\
        &\hspace{4cm}\cdot\left(\sum_{\substack{\sigma\in\thealtset\\ \sigma=\sigma_1\cdots \sigma_{m-1}\sigma_{m+1}\cdots\sigma_n}}\prod_{\substack{x=1\\x\neq m}}^{n(I^c)} (-1)^{\ell(\sigma_x)}\qpartition{\alpha_{i'_x,j'_x}-\sum_{i\in\mathcal{I}_x}\alpha_i}\right).\label{break into product 3}
    \end{align}
    Notice that by \Cref{break into product 3}, we have expressed $m_q(\hr_r,\alpha_I)$ as a product, with every term corresponding to some $\mathcal{I}_m$ contained in a single factor.
    By repeating this factorization for every remaining $i\in[n(I^c)]$, the weight $q$-multiplicity can be stated as
    \[m_q(\hr,\alpha_I)=\prod_{x=1}^{n(I^c)}\sum_{\substack{\sigma\in\mathcal{A}(\hr,\alpha_I) \\ \sigma=\sigma_x}}(-1)^{\ell(\sigma_x)}\qpartition{\alpha_{i'_x,j'_x}-\sum_{i\in\mathcal{I}_x}\alpha_i},\]
    which completes the proof.
\end{proof}

\end{document}
